% Current synthesis for Lectures 8--11. Included by synthesis.tex.

\Needspace{12\baselineskip}
\section{Lecture 8 --- No-Arbitrage and the Fundamental Theorems}
\label{syn:lecture8}
\lecturesummary{State-by-state payoff comparisons give dominance,
replication, and bounds. In a finite market, strictly positive state prices
are equivalent to no arbitrage. Completeness determines whether all claims
can be replicated and whether the pricing measure is unique.}

\subsection{Dominance, replication, and forward value}
In a frictionless market with short sales, the same payoff at the same date
must have the same price. If \(H_A\geq H_B\) in every state, then
\(V_t(H_A)\geq V_t(H_B)\): otherwise buy \(A\), sell \(B\), and invest
the initial receipt. The future payoff is nonnegative and the financed
receipt is positive. An arbitrage may equivalently be presented as a
zero-cost trade with a nonnegative, sometimes positive payoff.

A portfolio that reproduces \(H\) state by state gives its price by the
law of one price. For a stock paying continuous yield \(q\), constant
rates, and \(\tau=T-t\),
\[
\boxed{V_t(S_T-K)=S_t e^{-q\tau}-Ke^{-r\tau},\qquad
F(t,T)=S_t e^{(r-q)\tau}.}
\]
The first is the value of an existing forward with fixed delivery price
\(K\); the second is the delivery price making a new forward worth zero.
The prepaid share \(S_t e^{-q\tau}\) delivers one share at \(T\).

\subsection{Finite markets: state prices and martingale measures}
Take \(K\) terminal states, all with positive physical probability.
Let \(A\in\R^{K\times(m+1)}\) be the payoff matrix: rows are states and
columns are traded assets. Their current price vector is \(p\). The bank
account costs 1 and pays the gross return \(R>0\) in every state.

\begin{theorem}[Finite-state First Fundamental Theorem]
The market is arbitrage-free if and only if there exists a strictly
positive state-price vector \(\pi\) such that
\[
\boxed{A^\top\pi=p,\qquad \pi_k>0.}
\]
Since \(\sum_k\pi_k=1/R\), the numbers \(q_k=R\pi_k\) define an
equivalent martingale measure, and
\[
p_i=\frac1R\E^{\Qbb}[S_T^i].
\]
\end{theorem}
Equivalence means that no physically possible state is assigned zero
probability. A physical expected return is not itself a state price.
In a general continuous-time semimartingale market, the precise theorem
uses a suitable no-free-lunch condition and equivalent local martingale
measures; the elementary finite-state statement does not remove the need
for admissibility and integrability in continuous time.

\begin{theorem}[Finite-state Second Fundamental Theorem]
In an arbitrage-free finite market, the following are equivalent:
every payoff is replicable; \(\operatorname{rank}(A)=K\); and the
equivalent martingale measure is unique.
\end{theorem}
Adding an asset helps only when its payoff adds a missing direction.
Redundant traded assets can make holdings nonunique while leaving every
replication price unique. In a Brownian market, the corresponding dynamic
condition is that the traded volatility matrix spans all Brownian risks
(Lecture~7).

\subsection{Superhedging bounds and the price of unspanned risk}
For a claim \(h\), a subhedge pays at most \(h\), and a superhedge pays
at least \(h\). Their costs bound any compatible claim price. In a finite
arbitrage-free market, with \(\mathcal M\) the equivalent martingale
measures, linear-programming duality gives
\[
\begin{aligned}
\underline v(h)
&=\sup_{A\theta\leq h}p^\top\theta
=\inf_{\Qbb\in\mathcal M}\frac1R\E^{\Qbb}[h],\\
\overline v(h)
&=\inf_{A\theta\geq h}p^\top\theta
=\sup_{\Qbb\in\mathcal M}\frac1R\E^{\Qbb}[h].
\end{aligned}
\]
The extreme expectations can require limits of equivalent measures. When
a claim is added as a freely traded asset, its admissible prices are the
expectations generated by equivalent measures. Endpoints of the closed
hedging bounds need not themselves be admissible prices. A replicable
claim has a single price under every measure in \(\mathcal M\).

\paragraph{Executable quotes.}
With bid \(b_i\) and ask \(a_i\), dominance of \(H_B\) over \(H_A\)
requires \(b_A\leq a_B\). If the payoffs agree, both
\(b_A\leq a_B\) and \(b_B\leq a_A\) are necessary: the two quote
intervals must overlap. Unequal midpoints alone do not identify an
executable arbitrage, and these pairwise checks do not test all other
market constraints.

\Needspace{8\baselineskip}
\subsection*{Lecture 8 checklist}
\begin{itemize}[leftmargin=1.5em,itemsep=0.15em]
\item describe the purchase, sale, financing, and terminal payoff of an arbitrage;
\item distinguish a forward's current value from its delivery price;
\item solve \(A^\top\pi=p\), require strict positivity, and convert state prices to probabilities;
\item use a rank test to distinguish spanning from redundant instruments;
\item distinguish closed hedging bounds from admissible prices and executable bid--ask quotes.
\end{itemize}
\paragraph{Practice.} Core Exercises 8.1--8.4 cover arbitrage financing,
completeness, bounds, and executable quotes. Optional Exercise 7.11
extends the incomplete-market pricing argument.

\Needspace{12\baselineskip}
\section{Lecture 9 --- Vanilla Options in Practice}
\label{syn:lecture9}
\lecturesummary{Read the payoff and financing convention before using a
formula. Separate prices, implied volatilities, and realized variance.
Use Greeks as local exposures, and test quote consistency across strikes
and maturities. A hedging outcome depends on the timing of variance, not
only on an unweighted realized-volatility number.}

\subsection{Payoff, price, and profit are different quantities}
A call pays \((S_T-K)^+\), a put pays \((K-S_T)^+\), and a short position
reverses the payoff. Moneyness compares spot with strike. A straddle is a
call plus a put at the same strike; stock plus put creates a floor, and
stock minus call creates a cap. A profit calculation also includes the
initial premium and its financing. The terminal payoff alone is not P\&L.

For continuous yield \(q\), constant \(r\), and matched \(K,T\),
\[
\boxed{C_t-P_t=S_t e^{-q\tau}-Ke^{-r\tau}.}
\]
Dominance gives
\[
\begin{aligned}
(S_t e^{-q\tau}-Ke^{-r\tau})^+&\leq C_t\leq S_t e^{-q\tau},\\
(Ke^{-r\tau}-S_t e^{-q\tau})^+&\leq P_t\leq Ke^{-r\tau}.
\end{aligned}
\]
These bounds use discounted cash and prepaid stock. For a European option,
undiscounted intrinsic value is not a universal lower bound with arbitrary
carry or negative rates.

\subsection{Greeks, short maturity, and units}
The BS formula and Greeks are given in Lecture~5. Delta gives the share
holding for a replicating long claim; a seller holding the opposite claim
buys that replication. Gamma measures the change in delta. At fixed spot,
call delta tends to 0, \(1/2\), or 1 as expiry approaches according as
spot lies below, at, or above strike. The terminal kink has no classical
delta at equality.

At zero rates, no dividends, and ATM,
\[
\boxed{C_{\rm ATM}\sim\frac{S\sigma\sqrt\tau}{\sqrt{2\pi}},\quad
C_{\rm ATM}+P_{\rm ATM}\sim\frac{2S\sigma\sqrt\tau}{\sqrt{2\pi}},\quad
\Gamma_{\rm ATM}\sim\frac1{S\sigma\sqrt{2\pi\tau}}.}
\]
A fourfold maturity multiplies the leading ATM cost by two. Volatility is
annualized, with time in years; a daily return standard deviation \(d\)
corresponds to \(d\sqrt{252}\) under the trading-day convention. A vega
quoted per unit of decimal volatility must be multiplied by 0.01 for a
one-percentage-point change.

At fixed \(S,\sigma>0\) in BS with no dividends,
\[
\boxed{\mathrm{Vega}=S^2\sigma\tau\Gamma.}
\]
Consequently, same-expiry options do not independently span arbitrary
gamma and vega targets. Two hedge-option columns have determinant
\(S^2\sigma\Gamma_g\Gamma_h(\tau_h-\tau_g)\). Different maturities and
nonzero gammas give an invertible local sensitivity system; near
singularity can still produce unstable holdings. Positive payoff alone
does not imply positive vega: a bull call spread has changing curvature.

\subsection{Historical, realized, and implied volatility}
For observed log returns \(R_i\), historical annualized volatility is
\(\sqrt A\,s_R\), where \(A\) observations are made per year and
\(s_R\) is their sample standard deviation. The square-root rule assumes
uncorrelated returns with a common variance. Realized annualized variance
over a window \(T\) is
\[
\mathrm{RV}_{\rm ann}=\frac1T\sum_iR_i^2
\quad\longrightarrow\quad\frac1T\int_0^T\sigma_u^2\dd u
\]
for increasingly fine sampling of a continuous diffusion. Observation
noise changes this conclusion: for \(Y_i=\sigma W_{iT/n}+\varepsilon_i\)
with independent errors of variance \(\nu^2\),
\(\E\sum_i(\Delta Y_i)^2=\sigma^2T+2n\nu^2\).

Implied volatility inverts a market quote through the BS formula:
\[
C^{\rm mkt}(K,T)=C^{\rm BS}(S_0,K,T,r,\sigma_{\rm imp}(K,T)).
\]
It is a price in volatility units, not an instantaneous coefficient or a
historical forecast. Parity-consistent calls and puts with the same strike
and maturity have the same implied volatility. A downside skew compares
different strikes, rather than matched call and put contracts.

\subsection{The sign and timing of volatility-hedging error}
Assume no dividends, continuous costless rebalancing, and constant \(r\).
A call seller uses \(v(t,S_t)=C^{\rm BS}(S_t,K,T-t,r,\widehat\sigma)\)
and holds \(\Delta=v_s\). Under the actual dynamics, instantaneous
variance is \(\sigma_t^2\). With self-financing wealth \(X\), define
\(\Pi_t=X_t-v(t,S_t)\). Subtracting It\^o's formula from the hedge gain
and using the model PDE gives
\[
\dd\Pi_t=r\Pi_t\dd t+
\tfrac12(\widehat\sigma^2-\sigma_t^2)S_t^2\Gamma_t\dd t.
\]
Hence
\[
\boxed{\Pi_T=e^{rT}\Pi_0+
\frac12\int_0^Te^{r(T-u)}S_u^2\Gamma_u
(\widehat\sigma^2-\sigma_u^2)\dd u.}
\]
For a fairly funded initial hedge, \(\Pi_0=0\). The long-option convention
reverses the sign. For positive gamma, pointwise higher actual variance
hurts the seller, but a comparison of unweighted realized and quoted
volatility does not determine the weighted integral. Discrete hedging,
costs, and a different initial premium introduce additional effects.

\subsection{An arbitrage-consistent price surface}
At fixed expiry, a call is decreasing and convex in strike. With
\(D=e^{-r\tau}\),
\[
0\leq C(K_1)-C(K_2)\leq D(K_2-K_1),\qquad K_1<K_2.
\]
For unequal strikes \(K_1<K_2<K_3\), convexity requires
\[
C(K_2)\leq
\frac{K_3-K_2}{K_3-K_1}C(K_1)
+\frac{K_2-K_1}{K_3-K_1}C(K_3).
\]
The corresponding weighted butterfly has nonnegative terminal payoff.
Equal weights \((1,-2,1)\) apply only to equal strike spacings. Smooth
prices satisfy \(-D\leq C_K\leq0\) and \(C_{KK}\geq0\).

For the calendar condition, take no dividends, fix today's state, and
follow log-forward moneyness \(x\):
\[
K(T)=S_0e^{rT+x},\qquad c(T,x)=\frac{C(T,K(T))}{S_0}.
\]
If the discounted stock is a true martingale, conditional Jensen
applied to its positive-part payoff makes \(c(T,x)\) nondecreasing, so
\[
\boxed{c_T\big|_x=\frac{C_T+rKC_K}{S_0}\geq0.}
\]
In implied-volatility units this is nondecreasing total variance
\(w(T,x)=T\sigma_{\rm imp}(T,x)^2\) at fixed \(x\). Annualized implied
volatility may fall while total variance rises. A check at selected
strikes or maturities does not certify an entire interpolated surface.

\Needspace{8\baselineskip}
\subsection*{Lecture 9 checklist}
\begin{itemize}[leftmargin=1.5em,itemsep=0.15em]
\item distinguish payoff, intrinsic value, premium, and financed profit;
\item use parity and bounds before solving for an implied volatility;
\item track volatility and vega units and short-time square-root scaling;
\item interpret gamma-weighted hedging error with a stated seller or buyer convention;
\item check unequal-strike convexity and calendar behavior at fixed forward moneyness.
\end{itemize}
\paragraph{Practice.} Core Exercises 9.1--9.4 test sensitivity signs,
variance weighting, carry, and unequal-strike quote consistency.

\Needspace{12\baselineskip}
\section{Lecture 10 --- American Options and Optimal Stopping}
\label{syn:lecture10}
\lecturesummary{The payoff may be unchanged, but the holder now chooses
when it is received. Compare exercise with conditional continuation,
using only current information. Backward induction, the Snell envelope,
and the obstacle PDE describe the same decision problem. Numerical methods
must estimate an admissible exercise policy, not select dates with hindsight.}

\subsection{The pricing problem and the information constraint}
In the complete-market benchmark, with constant rate \(r\),
\[
\boxed{V_t^{\rm Am}=\esssup_{\tau\in\mathcal T_{t,T}}
\E^{\Qbb}[e^{-r(\tau-t)}g(S_\tau)\mid\F_t].}
\]
Here \(\mathcal T_{t,T}\) contains stopping times taking values in
\([t,T]\). The decision at \(u\) may use \(\F_u\), but not future prices.
The supremum of expected discounted payoffs is generally smaller than the
expectation of their pathwise maximum. American value is at least European
value because \(\tau=T\) is available, but this does not imply that early
exercise is optimal at a given state.

In a complete market under suitable assumptions, this optimal-stopping
value equals the seller's least initial capital for a hedge covering every
possible exercise date. Terminal replication alone is insufficient. In an
incomplete market, selecting one pricing measure does not by itself prove
equality with the least hedge cost across all exercise scenarios.

\subsection{Backward induction and the Snell envelope}
On dates \(t_0,\ldots,t_N\), let \(G_k=g(S_{t_k})\) and let \(B_k\) be
the bank account. Define the discounted value \(Z_k=V_k/B_k\). Then
\[
\boxed{Z_N=\frac{G_N}{B_N},\qquad
Z_k=\max\left\{\frac{G_k}{B_k},\E^{\Qbb}[Z_{k+1}\mid\F_{t_k}]\right\}.}
\]
The Snell envelope \(Z\) is the least supermartingale dominating
\(G_k/B_k\). Exercise is optimal when the exercise value equals the
envelope, using the first such contact as a standard optimal rule.
Its supermartingale decomposition is \(Z=M-A\), with increasing
compensator \(A\). In a complete market, a martingale representation of
\(M\) constructs a dominating self-financing hedge; the envelope itself
need not be self-financing wealth without consumption.

For a no-dividend binomial stock with multipliers \(d<u\), step
\(h\), and gross bank return \(R=e^{rh}\), no arbitrage requires
\(d<R<u\). The risk-neutral up probability and recursion are
\[
p=\frac{R-d}{u-d},\qquad
V_{k,j}=\max\left\{g(S_{k,j}),
R^{-1}[pV_{k+1,j+1}+(1-p)V_{k+1,j}]\right\}.
\]
At a continuation node, the one-step hedge uses
\[
\Delta=\frac{V_{\rm up}-V_{\rm down}}{S(u-d)},\qquad
b=\frac{uV_{\rm down}-dV_{\rm up}}{R(u-d)},
\]
where \(b\) is current cash. At exercise, the writer settles the exercise
payment. Keep annual rates, per-step gross returns, and the exercise grid
consistent.

\subsection{The obstacle PDE and the free boundary}
With continuous yield \(q\), define the pricing generator
\[
\mathcal L u=(r-q)su_s+\tfrac12\sigma^2s^2u_{ss}.
\]
The value satisfies \(u\geq g\) and
\(u_t+\mathcal L u-ru\leq0\). In the continuation region, waiting is
optimal and the latter inequality is an equality. Formally,
\[
\boxed{\max\{g-u,\ u_t+\mathcal L u-ru\}=0,\qquad u(T,s)=g(s).}
\]
For a regular American put boundary \(b(t)\),
\[
u(t,b(t))=K-b(t),\qquad u_s(t,b(t))=-1.
\]
These are value matching and smooth fit. The boundary is unknown and must
be found together with the continuation value. For the no-dividend put
with \(r\geq0\), the spatial conditions include \(u(t,0)=K\) and
\(u(t,s)\to0\) as \(s\to\infty\).

For a no-dividend call and \(r\geq0\),
\(C^{\rm Eur}\geq s-Ke^{-r\tau}\geq s-K\): early exercise brings no
advantage, so American and European call values agree. Dividends or
negative rates can change the conclusion. Inside smooth in-the-money
exercise regions, necessary local tests with yield \(q\) are
\[
\boxed{qs\leq rK\quad\text{for a put},\qquad
qs\geq rK\quad\text{for a call}.}
\]
They follow by inserting the exercise payoff into
\((\mathcal L-r)g\leq0\); they are not sufficient global exercise rules.

\subsection{Three numerical routes and what each approximates}
\paragraph{Trees.}
Compute terminal payoffs and recurse backward, keeping both values and
exercise decisions. A time-discrete exercise grid approximates continuous
exercise; the stock-tree approximation is another source of error.

\paragraph{Longstaff--Schwartz simulation.}
On simulated paths and a chosen exercise grid:
\begin{enumerate}[leftmargin=1.5em,itemsep=0.15em]
\item work backward from the final exercise date;
\item on relevant in-the-money training paths, regress the realized later
cash flows, discounted from their already selected exercise dates, on
functions of the current state;
\item exercise when intrinsic value exceeds fitted conditional continuation;
\item value the resulting fixed, adapted policy on independent paths,
using actual realized cash flows selected by that policy.
\end{enumerate}
The fitted conditional mean is not a guaranteed payment. Taking the maximum
of intrinsic value and each path's realized future payment introduces
look-ahead. Independent evaluation separates policy fitting from selection
bias: its true policy expectation is at most the optimal value, while a
finite-sample estimate still has statistical error. Basis choice, sample
size, state variables, and the exercise grid all matter.

\paragraph{Finite differences.}
Discretize the backward continuation equation, then impose the obstacle.
For a centered explicit step on \(s_j=j\Delta s\),
\[
\widetilde U_j=A_jU_{j-1}+B_jU_j+C_jU_{j+1},\qquad
U_j^{\rm new}=\max\{g(s_j),\widetilde U_j\},
\]
where
\[
\begin{aligned}
A_j&=\tfrac12h[\sigma^2j^2-(r-q)j],\\
B_j&=1-h(\sigma^2j^2+r),\\
C_j&=\tfrac12h[\sigma^2j^2+(r-q)j].
\end{aligned}
\]
Nonnegative coefficients support a monotone continuation update. For
\(r\geq0\) and interior indices up to \(J\),
\(h\leq1/(\sigma^2J^2+r)\) ensures \(B_j\geq0\); separately,
\(\sigma^2j\geq|r-q|\) is needed for \(A_j,C_j\geq0\).
An obstacle projection does not repair negative continuation weights.
Check the spatial scheme, time step, truncation boundary, and exercise
condition separately.

\paragraph{Variation: settlement after exercise.}
If \(g(S_\tau)\) is fixed at exercise but paid a constant \(\delta\)
years later, its time-\(\tau\) reward is \(e^{-r\delta}g(S_\tau)\).
Thus the entire value is multiplied by \(e^{-r\delta}\) and the
optimal exercise region is unchanged. State-dependent delays or a payment
recomputed from the later stock price require a different exercise reward.

\Needspace{8\baselineskip}
\subsection*{Lecture 10 checklist}
\begin{itemize}[leftmargin=1.5em,itemsep=0.15em]
\item state the allowed stopping times and the information available at exercise;
\item distinguish immediate payoff, conditional continuation, and a realized future cash flow;
\item write backward induction, the Snell envelope, and the obstacle with consistent discounting;
\item state the rate and dividend assumptions behind early-exercise claims;
\item assess exercise-policy error, sampling error, and PDE or tree discretization separately.
\end{itemize}
\paragraph{Practice.} Core Exercises 10.1--10.4 cover the writer's
hedge, delayed settlement, dividend-aware exercise, and independent
evaluation of a learned stopping policy.

\Needspace{12\baselineskip}
\section{Lecture 11 --- Beyond Black--Scholes: Dupire and Heston}
\label{syn:lecture11}
\lecturesummary{A vanilla surface reveals terminal pricing distributions.
Local volatility constructs one diffusion consistent with suitable surface
data. Heston introduces random variance, an additional risk premium, and
an option-based hedge. Calibration and simulation must be checked separately
from the model's assumptions about future dynamics.}

\subsection{Model coefficients and contract-price derivatives}
Assume constant rate \(r\), no dividends unless stated otherwise, European
payoffs, and a selected pricing measure \(\Qbb\). Distinguish
\[
u(t,s;T,K)\quad\text{from}\quad C(T,K)=u(0,S_0;T,K).
\]
The backward PDE varies valuation state \((t,s)\) for a fixed contract;
the forward price equation varies maturity and strike \((T,K)\) with
today's state fixed. In particular, \(C_T\) is not the calendar-time theta
\(u_t\), and \(C_K\) is not the stock delta \(u_s\).

\begin{center}
\begin{tabular}{L{.20\linewidth}L{.33\linewidth}L{.35\linewidth}}
\toprule Model & State and Brownian risks & Volatility assumption\\\midrule
Black--Scholes & \((t,S_t)\), one risk & Constant instantaneous volatility\\[.3em]
Local volatility & \((t,S_t)\), one risk & Deterministic function \(\sigma_{\rm loc}(t,s)\)\\[.3em]
Heston & \((t,S_t,V_t)\), two risks in the nondegenerate case & Random variance with specified pricing dynamics\\
\bottomrule
\end{tabular}
\end{center}
In the local model,
\[
\dd S_t=rS_t\dd t+\sigma_{\rm loc}(t,S_t)S_t\dd W_t^{\Qbb},
\]
and the backward pricing PDE remains
\[
u_t+rsu_s+\tfrac12\sigma_{\rm loc}^2(t,s)s^2u_{ss}-ru=0.
\]
A spot-dependent diffusion coefficient does not add an independent factor.

\subsection{Terminal densities and Dupire's forward equation}
Assume a smooth density \(p(T,x)\), with sufficient regularity to
interchange derivatives and integrals. Then
\[
C(T,K)=e^{-rT}\int_K^\infty(x-K)p(T,x)\dd x.
\]
Leibniz's rule first removes the payoff's strike dependence; the lower
boundary vanishes because \(x-K=0\) there. Differentiating again gives
\[
\boxed{C_K=-e^{-rT}\Qbb(S_T>K),\qquad C_{KK}=e^{-rT}p(T,K).}
\]
A tight call spread reveals a discounted tail probability, and call
curvature reveals the risk-neutral density. This is pricing information,
not an estimate of historical frequencies. With terminal atoms, curvature
is a measure and the one-sided strike derivatives differ at the atom;
the classical smooth-density quotient below is not justified there.

Set \(a(T,x)=\sigma_{\rm loc}^2(T,x)x^2\). It\^o's formula for a compactly
supported test function, followed by one drift and two diffusion
integrations by parts, yields the forward density equation
\[
p_T=-\partial_x(rxp)+\tfrac12\partial_{xx}(ap).
\]
Differentiating the discounted call expectation at fixed strike,
\[
C_T=-rC+e^{-rT}(I_{\rm drift}+I_{\rm diff}).
\]
Integration by parts gives
\[
\begin{aligned}
I_{\rm drift}&=r\int_K^\infty xp(T,x)\dd x,\\
I_{\rm diff}&=\tfrac12a(T,K)p(T,K),\\
e^{-rT}\int_K^\infty xp(T,x)\dd x&=C-KC_K.
\end{aligned}
\]
Besides smoothness, assume the tail boundary terms
\((x-K)rxp\), \((x-K)\partial_x(ap)\), and \(ap\) vanish at infinity.
The first lower-boundary terms vanish at \(x=K\); the remaining diffusion
term evaluates the coefficient at that strike. Consequently,
\[
C_T=-rC+r(C-KC_K)+\tfrac12\sigma_{\rm loc}^2K^2C_{KK},
\]
so
\[
\boxed{C_T+rKC_K=\tfrac12\sigma_{\rm loc}^2(T,K)K^2C_{KK},\qquad
\sigma_{\rm loc}^2(T,K)=\frac{C_T+rKC_K}{\tfrac12K^2C_{KK}}.}
\]
The quotient is used where \(K>0\) and \(C_{KK}>0\). A flat BS surface
recovers its constant variance. The value at \((T,K)\) is a local model
coefficient at time \(T\), spot \(K\), not the implied volatility of the
contract \((T,K)\).

\paragraph{Variation: continuous dividends.}
For stock drift \(r-q\) with constant yield \(q\), discounting is still
at \(r\). The price-level terms no longer cancel completely:
\[
\boxed{\sigma_{\rm loc}^2(T,K)
=\frac{C_T+qC+(r-q)KC_K}{\tfrac12K^2C_{KK}}.}
\]
Using the no-dividend formula with nonzero yield loses both the correct
carry coefficient and the term \(qC\).

\subsection{Calibration of local volatility is a numerical problem}
A smooth, arbitrage-consistent surface is needed before differentiating.
At equally spaced strikes,
\[
C_{KK}\approx\frac{C(K+h)-2C(K)+C(K-h)}{h^2}.
\]
Independent quote errors of standard deviation \(\varepsilon\) give a
curvature-error standard deviation \(\sqrt6\,\varepsilon/h^2\).
Small spacing amplifies noise, and small curvature makes the Dupire ratio
unstable. Clipping a negative variance does not repair inconsistent prices.
Check bounds, slopes, convexity, carry-adjusted maturity behavior,
interpolation, and extrapolation.

For the lecture's five-quote exercise,
\(C_T\simeq4\), \(C_K\simeq-0.49\), \(C_{KK}\simeq0.02\),
\(K=100\), and \(r=0.05\). Thus
\[
\sigma_{\rm loc}^2\simeq\frac{4-2.45}{100}=0.0155,
\qquad \sigma_{\rm loc}\simeq12.45\%.
\]
The carry correction changes the answer substantially.

A full vanilla surface determines the marginal law of \(S_T\) at each
maturity. It does not determine the full joint path law. Local volatility
supplies one suitable diffusion consistent with those marginals; other
dynamics can imply different average, maximum, or barrier-event values.
A match of all available vanillas is therefore not a validation of every
exotic price or hedge.

\subsection{Heston dynamics, the variance mean, and the boundary}
Under the chosen pricing measure,
\[
\boxed{\begin{aligned}
\frac{\dd S_t}{S_t}&=r\dd t+\sqrt{V_t}\dd W_t^S,\\
\dd V_t&=\kappa(\theta-V_t)\dd t+\eta\sqrt{V_t}\dd W_t^V,
\qquad\dd\langle W^S,W^V\rangle_t=\rho\dd t.
\end{aligned}}
\]
Here \(V_0\) is current variance, \(\theta\) its long-run level,
\(\kappa\) mean-reversion speed, \(\eta\) volatility of variance, and
\(\rho\) stock--variance shock correlation. The state is \((S_t,V_t)\),
so equal spots need not imply equal option prices.

The conditional variance mean is
\[
\boxed{\E^{\Qbb}[V_T\mid V_t=v]
=\theta+(v-\theta)e^{-\kappa(T-t)}.}
\]
It is not an implied variance, and its square root is not generally
\(\E[\sqrt{V_T}\mid V_t]\). With positive CIR parameters, variance is
nonnegative. Starting at \(V_0>0\), the Feller condition
\(2\kappa\theta\geq\eta^2\) makes zero unattainable for the exact
continuous process. It does not make an Euler step nonnegative.

\paragraph{Expected average variance and the ATM guide.}
Integrating the conditional mean gives the exact identity
\[
\boxed{w_t(\tau):=\frac1\tau\E^{\Qbb}\!\left[\int_t^{t+\tau}V_u\dd u
\middle|\F_t\right]
=\theta+(V_t-\theta)\frac{1-e^{-\kappa\tau}}{\kappa\tau}.}
\]
The limits are \(V_t\) for \(\tau\downarrow0\) and \(\theta\) for
\(\tau\to\infty\). Larger \(\kappa\) reduces the persistence of
\(V_t-\theta\). A linear payment of average variance uses this mean,
with discounting applied separately. In contrast,
\(\Sigma_{\rm ATM}^2(\tau)\approx w_t(\tau)\) is a deterministic-variance
guide, exact when \(\eta=0\), not an identity for stochastic variance.
The long-maturity limit of the mean does not establish
\(\Sigma_{\rm ATM}^2\to\theta\). When \(V_0=\theta=0.04\), the mean
is flat at 0.04, while the lecture's calculated ATM volatilities are about
19.30\% at three months and 18.74\% at two years.

\subsection{The variance risk premium is part of the specification}
Use independent Brownian coordinates \(Z_1,Z_2\) and write
\(\dd W^V=\rho\dd Z_1+\sqrt{1-\rho^2}\dd Z_2\), with stock driven by
\(Z_1\). Under \(\Pbb\), let the stock drift be \(\mu\) and the
variance drift \(b_{\Pbb}(v)\). A valid Girsanov change with
\(\dd Z_j^{\Qbb}=\dd Z_j^{\Pbb}+\lambda_j\dd t\) gives, on \(v>0\),
\[
\lambda_1=\frac{\mu-r}{\sqrt v},\qquad
b_{\Qbb}=b_{\Pbb}-\pi,
\]
\[
\pi=\eta\sqrt v\bigl(\rho\lambda_1+
\sqrt{1-\rho^2}\lambda_2\bigr).
\]
The stock condition fixes \(\lambda_1\) but leaves the independent
variance-risk shift \(\lambda_2\) free. Stock-and-cash no-arbitrage alone
does not select \(b_{\Qbb}\). Under an admissible affine choice
\(\pi(v)=\lambda v\),
\[
\boxed{\kappa_{\Qbb}=\kappa_{\Pbb}+\lambda,\qquad
\theta_{\Qbb}=\frac{\kappa_{\Pbb}\theta_{\Pbb}}
{\kappa_{\Pbb}+\lambda},\qquad\kappa_{\Pbb}+\lambda>0.}
\]
This convention preserves \(\kappa\theta\). Equivalent drift changes
preserve the fixed processes' quadratic variations and covariations;
they do not freely change \(\eta\) or \(\rho\). Matching the drift
algebra still requires a true, equivalent density martingale, particularly
when shifts can become singular near \(v=0\).
For the lecture's example, \(\kappa_{\Pbb}=1.6\),
\(\theta_{\Pbb}=0.05\), and \(\lambda=0.4\) give
\(\kappa_{\Qbb}=2\), \(\theta_{\Qbb}=0.04\), with common product 0.08.

\subsection{Pricing PDE and stock-plus-option replication}
For \(u(t,s,v)\), quadratic covariation adds the mixed It\^o term
\(\rho\eta svu_{sv}\). The two symmetric cross terms combine, so its
coefficient has no extra factor \(1/2\). Under the selected pricing
dynamics, the European PDE is
\[
\boxed{\begin{aligned}
0={}&u_t+rsu_s+\kappa(\theta-v)u_v+\tfrac12vs^2u_{ss}\\
&+\rho\eta svu_{sv}+\tfrac12\eta^2vu_{vv}-ru,
\qquad u(T,s,v)=h(s).
\end{aligned}}
\]
Holding \(u_s\) shares removes the direct stock term but leaves variance
exposure. With \(|\rho|<1\), no choice of stock holding removes the
independent variance shock when \(\eta\sqrt v\,u_v\ne0\).

Let \(g(t,s,v)\) denote a suitable traded hedge option, priced with
the same dynamics. For \(v,\eta>0\), \(|\rho|<1\), and \(g_v\ne0\),
matching both Brownian exposures gives
\[
\boxed{\alpha=\frac{u_v}{g_v},\qquad
\Delta=u_s-\alpha g_s,\qquad b=u-\Delta s-\alpha g.}
\]
To hedge one short claim, hold \(\alpha\) hedge options, \(\Delta\)
shares, and cash amount \(b\), and rebalance self-financingly. Replication
also needs an available, nondegenerate hedge instrument over the horizon,
regularity, and admissible trading. Small \(g_v\) can make positions
unstable. Variance derivatives are with respect to the model state, not
a contract-specific BS implied-volatility input. Two traded options must
infer the same variance drift from their respective pricing PDEs if they
are to be consistent with one pricing model.

\paragraph{The drift condition, without dividing by the target sensitivity.}
Let \(\mathcal A^{\Pbb}f=f_t+\mu sf_s+b_{\Pbb}(v)f_v+
vs^2f_{ss}/2+\rho\eta svf_{sv}+\eta^2vf_{vv}/2\).
Matching the self-financing portfolio drift along the historical paths yields
\[
\mathcal A^{\Pbb}u-ru-(\mu-r)su_s
=\frac{u_v}{g_v}
\bigl[\mathcal A^{\Pbb}g-rg-(\mu-r)sg_s\bigr]=\pi u_v.
\]
Thus the claim's PDE uses \(b_{\Pbb}-\pi=b_{\Qbb}\), and the
historical stock drift cancels. The identity remains valid where \(u_v=0\).
One liquid option with nonzero variance sensitivity can span the missing
Brownian exposure in a specified model. One observed option price does
not identify the full function \(g(t,s,v)\), its derivatives, or the premium.
Calibration of a set of option quotes selects the model's pricing drift;
estimating \(\pi=b_{\Pbb}-b_{\Qbb}\) separately additionally requires
historical variance-drift information. Prices and hedge ratios use the
chosen pricing dynamics without that separate decomposition.

\subsection{Calibration, parameter effects, and identification}
For vanilla quotes \(C_i^{\rm mkt}\), a price-based objective is
\[
\min_{\vartheta}\sum_i w_i
\bigl[C^{\rm Heston}(K_i,T_i;\vartheta)-C_i^{\rm mkt}\bigr]^2,
\qquad\vartheta=(V_0,\kappa,\theta,\eta,\rho).
\]
State instrument choices, weights, bid--ask treatment, and constraints.
Check several starting points and parameter stability. These are pricing
parameters under \(\Qbb\), rather than historical estimates under
\(\Pbb\). One ATM quote cannot regularly identify five parameters: its
Jacobian has rank at most one. Even many quotes need a full-column-rank,
well-conditioned sensitivity matrix for stable local identification;
this does not ensure a unique global fit.

\paragraph{State, pricing parameters, and historical parameters.}
\((S_t,V_t)\) is the current state; \(V_t\) is inferred rather than directly
observed. Within a fixed equivalent-measure Heston specification, the
diffusion inputs \(\eta,\rho\) are shared by \(\Pbb\) and \(\Qbb\),
whereas the variance-drift parameters \(\kappa,\theta\) may differ.
Option calibration estimates pricing inputs; time-series estimation addresses
historical dynamics. A mismatch of fitted correlations is not freely
explained by changing measure in the same model.

\begin{center}
\begin{tabular}{cL{.76\linewidth}}
\toprule Parameter & Effect to examine in the lecture's controlled comparisons\\\midrule
\(\rho\) & More negative correlation strengthens downside skew in the experiment.\\[.3em]
\(\eta\) & Changes wings, curvature, and ATM level together.\\[.3em]
\(V_0\) & Has a stronger short-maturity influence in the experiment.\\[.3em]
\(\theta\) & Has a stronger long-maturity influence in the experiment.\\[.3em]
\(\kappa\) & Changes variance persistence and distribution, even when \(V_0=\theta\) keeps the mean variance constant.\\
\bottomrule
\end{tabular}
\end{center}
These are model comparisons around the illustrative pricing baseline
\[
(V_0,\kappa,\theta,\eta,\rho)=(0.04,2,0.04,0.4,-0.6),
\qquad S_0=100,\quad r=0,
\]
at 0.25 and 2 years. They are not universal
comparative statics or market observations. A linear average-variance
payment depends only on the conditional variance mean, while nonlinear
option prices depend on richer distributional and dependence information.

\paragraph{Skew and maturity: a local approximation.}
For log-forward moneyness \(k=\log(K/(S_0e^{rT}))\) and the flat-mean
baseline \(V_0=\theta\), the leading small-\(\eta\) guide is
\[
\partial_k\Sigma^2(0,T)\approx\frac{\rho\eta}{\kappa T}
\left[1-\frac{1-e^{-\kappa T}}{\kappa T}\right].
\]
Its short-maturity volatility slope is \(\rho\eta/(4\sqrt{V_0})\).
The sign follows \(\rho\), its size scales with \(\eta\), and mean
reversion damps the longer-dated skew within the approximation. This
simple finite-maturity formula assumes \(V_0=\theta\); use the full
pricer for general parameters and calibration, rather than treating these
guides as exact comparative statics.

\paragraph{Parameter risk survives a state hedge.}
The stock-plus-option hedge cancels movements in \(S,V\) under fixed model
parameters. A refit of a parameter \(a\), with the state held fixed,
instead produces approximately
\((\alpha\,\partial_a g-\partial_a u)\Delta a\) for the hedged seller.
This need not vanish. A small calibration residual today and a good hedge
under tomorrow's dynamics are separate checks.

\subsection{Simulation: covariance, truncation, and separate errors}
For step \(h\), draw fresh independent \(Z_1,Z_2\sim N(0,1)\), and use
\[
\Delta W^S=\sqrt h Z_1,\qquad
\Delta W^V=\sqrt h\bigl(\rho Z_1+\sqrt{1-\rho^2}Z_2\bigr).
\]
Both variances are \(h\) and covariance is \(\rho h\). Check the whole
covariance matrix, not correlation alone. A full-truncation variance step
and log-stock step are
\[
\boxed{\begin{aligned}
V_{n+1}&=V_n+\kappa(\theta-V_n^+)h
+\eta\sqrt{V_n^+}\,\Delta W_n^V,\\
S_{n+1}&=S_n\exp\bigl[(r-V_n^+/2)h+\sqrt{V_n^+}\,\Delta W_n^S\bigr],
\qquad V_n^+=\max(V_n,0).
\end{aligned}}
\]
Keep the raw \(V_n\) in the first term; truncate inside coefficients.
The raw variance may be negative, while the stock stays positive.
Clamping the entire next variance defines a different scheme. Both
schemes require numerical validation rather than an appeal to the
continuous Feller condition.

For \(M\) independent payoff samples, the discounted estimate and its
sampling standard error are
\[
\widehat V_0=e^{-rT}\overline H_M,\qquad
\widehat{\rm SE}=e^{-rT}\frac{s_H}{\sqrt M}.
\]
Increasing \(M\) reduces sampling error. Refine \(h\) to assess time bias,
and check against an independent vanilla price and limiting cases.
When \(\eta=0\), variance is deterministic but need not be constant;
constant-volatility BS is recovered if also \(V_0=\theta\).
Neither more paths nor a narrow Monte Carlo interval resolves model risk
or poor calibration identification.

\Needspace{8\baselineskip}
\subsection*{Lecture 11 checklist}
\begin{itemize}[leftmargin=1.5em,itemsep=0.15em]
\item distinguish implied volatility, local volatility, and state variance;
\item derive strike probabilities and Dupire with consistent discounting, carry, and boundary terms;
\item explain why terminal marginals alone do not determine path-dependent prices;
\item specify the variance pricing drift before writing the Heston PDE;
\item identify the mixed derivative and the option exposure needed to span variance risk;
\item distinguish an expected variance from an implied variance, and a calibrated pricing drift from an estimated risk premium;
\item distinguish fit, parameter identification, sampling error, and discretization bias.
\end{itemize}
\paragraph{Practice.} Core Exercises 11.1--11.4 cover dividend-adjusted
Dupire, an average-variance payment, option-hedge consistency, and
simulation scaling. Optional Exercises 11.5--11.6 extend joint dynamics
and the variance risk premium.
